\providecommand{\Needspace}[1]{}
\Needspace{8\baselineskip}
\item \textbf{Clasificación de estrellas según su brillo aparente}.
\nopagebreak[4]

\begin{tcolorbox}[colback=blue!2!white,colframe=blue!55!black,title={Enunciado},breakable]
En astronomía, el brillo aparente de una estrella puede estimarse a partir de su luminosidad y su distancia al observador. Una forma simplificada de calcular este brillo es mediante la expresión:
\[
F = \frac{L}{4\pi d^2},
\]
donde $F$ representa el brillo observado, $L$ la luminosidad de la estrella y $d$ la distancia a la estrella.

Considere los siguientes arreglos con información de distintas estrellas:
\[
L = [120,\ 80,\ 15,\ 5,\ 200,\ 40],
\]
\[
d = [3,\ 5,\ 8,\ 10,\ 6,\ 12].
\]

Cada elemento de ambos arreglos corresponde a una misma estrella. Es decir, la luminosidad \texttt{L[i]} y la distancia \texttt{d[i]} pertenecen a la estrella ubicada en el índice \texttt{i}.

Escriba un programa en Python que realice lo siguiente:

\begin{itemize}
    \item[(a)] Defina los arreglos de luminosidad $L$ y distancia $d$ usando \texttt{np.array}.

    \item[(b)] Cree una lista vacía donde se almacenarán los valores de brillo $F$ calculados para cada estrella.

    \item[(c)] Recorra todos los índices de los arreglos usando un ciclo \texttt{for}. En cada iteración, acceda a \texttt{L[i]} y \texttt{d[i]} para calcular el brillo de la estrella correspondiente y agréguelo a la lista, por ejemplo con \texttt{lista\_vacia.append(calculo\_F)}:
    \[
    F_i = \frac{L_i}{4\pi d_i^2}.
    \]

    \item[(d)] Para cada estrella, clasifique su brillo usando condicionales según el siguiente criterio:
    \[
    F > 1 \quad \Longrightarrow \quad \text{Muy brillante},
    \]
    \[
    0.1 < F \leq 1 \quad \Longrightarrow \quad \text{Brillo medio},
    \]
    \[
    F \leq 0.1 \quad \Longrightarrow \quad \text{Débil}.
    \]

    \item[(e)] Muestre en pantalla, para cada estrella, su índice, su luminosidad, su distancia, su brillo calculado y su clasificación.

    \item[(f)] Calcule e imprima el brillo promedio de todas las estrellas.

    \item[(g)] Realice un gráfico de dispersión usando \texttt{plt.scatter}, donde el eje horizontal represente la distancia $d$ y el eje vertical represente el brillo aparente $F$. Agregue título, nombres de ejes, grilla y una presentación limpia de la figura.
\end{itemize}

El objetivo del ejercicio es practicar el recorrido de arreglos mediante índices, usando expresiones como \texttt{L[i]} y \texttt{d[i]}, junto con condicionales para clasificar datos físicos y visualizar los resultados mediante un gráfico de dispersión.
\end{tcolorbox}

\begin{solution}
Una posible solución consiste en recorrer los arreglos usando sus índices. En cada iteración se calcula el brillo de la estrella correspondiente, se guarda el resultado en una lista y se clasifica usando condicionales.

\begin{pythonjupyter}
import numpy as np
import matplotlib.pyplot as plt

# Arreglos de datos
L = np.array([120, 80, 15, 5, 200, 40])
d = np.array([3, 5, 8, 10, 6, 12])

# Lista vacía para guardar los brillos
brillos = []

# Recorrido usando índices
for i in range(len(L)):

    # Cálculo del brillo aparente
    F = L[i] / (4 * np.pi * d[i]**2)

    # Guardar el brillo calculado
    brillos.append(F)

    # Clasificación usando condicionales
    if F > 1:
        clasificacion = "Muy brillante"
    elif F > 0.1:
        clasificacion = "Brillo medio"
    else:
        clasificacion = "Débil"

    # Mostrar resultados de cada estrella
    print("=" * 40)
    print("Estrella índice:", i)
    print("Luminosidad:", L[i])
    print("Distancia:", d[i])
    print("Brillo:", F)
    print("Clasificación:", clasificacion)

# Convertir la lista de brillos a arreglo de NumPy
brillos = np.array(brillos)

# Brillo promedio
brillo_promedio = np.mean(brillos)

print("=" * 40)
print("Brillo promedio:", brillo_promedio)

# Gráfico de dispersión
plt.figure(figsize=(7, 5))
plt.scatter(d, brillos)

plt.title("Brillo aparente de estrellas según su distancia")
plt.xlabel("Distancia d")
plt.ylabel("Brillo aparente F")
plt.grid(True)

plt.show()
\end{pythonjupyter}
\end{solution}
